\hypertarget{deploying-infrastructure-as-code-with-terraform}{%
\paragraph{Deploying Infrastructure as Code with
Terraform}}

Instead of manually configuring each service on the AWS Console, the
NewsRAG project uses Terraform to automate the provisioning of the
entire infrastructure (Infrastructure as Code - IaC).

\textbf{1. Terraform Structure}

The \texttt{main.tf} file (over 370 lines) manages all resources: -
\textbf{Networking}: VPC (\texttt{10.0.0.0/16}), 2 Public Subnets
(multi-AZ), Internet Gateway, Route Tables. - \textbf{Security Groups}:
Access control following the Principle of Least Privilege --- e.g., RDS
only accepts traffic from the ECS Security Group. - \textbf{Database}:
RDS Aurora PostgreSQL Serverless cluster. - \textbf{Compute}: ECS
Cluster, ECR Repository, and Task Definitions for Crawler, ETL,
Vectorize running on Fargate. - \textbf{Scheduling}: EventBridge Rules
to trigger ECS Tasks via cron schedules (01:00, 02:00, 03:00 UTC).

\textbf{2. Deploying to AWS}

To deploy the system, you only need to configure the AWS CLI and run
these commands:

\begin{Shaded}
\begin{Highlighting}[]
\CommentTok{\# Initialize Terraform provider}
\ExtensionTok{terraform}\NormalTok{ init}

\CommentTok{\# Preview the resources to be created}
\ExtensionTok{terraform}\NormalTok{ plan}

\CommentTok{\# Execute resource creation on AWS}
\ExtensionTok{terraform}\NormalTok{ apply }\AttributeTok{{-}auto{-}approve}
\end{Highlighting}
\end{Shaded}

\textbf{3. Deployment Automation (CI/CD)}

The \texttt{deploy.sh} script is provided to automatically build Docker
images, tag and push them to Amazon ECR, and then update the ECS service
to force a deployment of the latest version:

\begin{Shaded}
\begin{Highlighting}[]
\CommentTok{\#!/bin/bash}
\CommentTok{\# deploy.sh}

\CommentTok{\# 1. Authenticate with ECR}
\ExtensionTok{aws}\NormalTok{ ecr get{-}login{-}password }\AttributeTok{{-}{-}region}\NormalTok{ ap{-}southeast{-}2 }\KeywordTok{|} \ExtensionTok{docker}\NormalTok{ login }\AttributeTok{{-}{-}username}\NormalTok{ AWS }\AttributeTok{{-}{-}password{-}stdin} \VariableTok{$ACCOUNT\_ID}\NormalTok{.dkr.ecr.ap{-}southeast{-}2.amazonaws.com}

\CommentTok{\# 2. Build Docker images}
\ExtensionTok{docker}\NormalTok{ build }\AttributeTok{{-}t}\NormalTok{ newsrag{-}crawler }\AttributeTok{{-}f}\NormalTok{ crawler/Dockerfile .}
\ExtensionTok{docker}\NormalTok{ build }\AttributeTok{{-}t}\NormalTok{ newsrag{-}etl }\AttributeTok{{-}f}\NormalTok{ etl/Dockerfile .}

\CommentTok{\# 3. Tag and Push to ECR}
\ExtensionTok{docker}\NormalTok{ tag newsrag{-}crawler:latest }\VariableTok{$ACCOUNT\_ID}\NormalTok{.dkr.ecr.ap{-}southeast{-}2.amazonaws.com/newsrag{-}api:crawler{-}latest}
\ExtensionTok{docker}\NormalTok{ push }\VariableTok{$ACCOUNT\_ID}\NormalTok{.dkr.ecr.ap{-}southeast{-}2.amazonaws.com/newsrag{-}api:crawler{-}latest}
\end{Highlighting}
\end{Shaded}

\textbf{4. Resource Cleanup}

When the system is not in use, to avoid incurring costs, you can delete
all resources with a single command:

\begin{Shaded}
\begin{Highlighting}[]
\ExtensionTok{terraform}\NormalTok{ destroy}
\end{Highlighting}
\end{Shaded}

\begin{quote}
\textbf{Note}: This command will delete all VPC, ECS, RDS, and
EventBridge resources defined in \texttt{main.tf}. Database data will be
lost if you haven't configured a snapshot.
\end{quote}
